\section*{Hoofdstuk \ref{ch:b2:orbital-forcing} --- Orbitale forcering en de ijstijden}

\begin{solution}{exo:b2:orbital-forcing:1}
(a) $1.0167/0.9833 = 1.0340$: het aphelium ligt \qty{3.4}{\percent} verder dan het perihelium.
(b) $a(1 + e) - a(1 - e) = 2ea = 0.0334\,a$, \qty{3.34}{\percent} van $a$, dicht bij de 3.4 procent
die gewoonlijk wordt opgegeven. (c) Het zonlicht neemt af met het omgekeerde kwadraat van de afstand: $1.0340^2 = 1.069$,
ongeveer \qty{7}{\percent} meer zonlicht in het perihelium.
\end{solution}

\begin{solution}{exo:b2:orbital-forcing:2}
(a) Vandaag $0.017236 \times \sin\ang{101.37} = 0.017236 \times 0.9804 = 0.01690$; \num{10000} jaar geleden
$0.019971 \times \sin\ang{293.86} = 0.019971 \times (-0.9145) = -0.01826$; \num{6000} jaar geleden
$0.019249 \times \sin\ang{359.99} \approx 0$, de waarden van de tabellen. (b) De lengte van de zon in het
perihelium is $\varpi + \ang{180}$: \ang{281.37} vandaag, \ang{473.86} $-$ \ang{360} $=$ \ang{113.86}
\num{10000} jaar geleden, \ang{179.99} \num{6000} jaar geleden. (c) \num{10000} jaar geleden viel het perihelium
\ang{24} na de junizonnewende, van de drie het dichtst erbij, in de noordelijke zomer;
\num{6000} jaar geleden viel het bij de septemberequinox, waar de \omterm{def:b2:orbital-forcing:elements}{klimatologische precessie} nul is.
\end{solution}

\begin{solution}{exo:b2:orbital-forcing:3}
(a) De frequenties tellen op: $1/25.77 + 1/112 = 0.03880 + 0.00893 = \qty{0.04773}{kyr^{-1}}$, een periode
van $1/0.04773 = 21.0$ duizend jaar. (b) Het resultaat valt binnen het bereik van de hoofdperioden,
19.0 tot 23.7 duizend jaar. Eén getal kan ze niet reproduceren: het draaien van de baan
is zelf de som van verscheidene bewegingen met verschillende frequenties, die elk met de
precessie van de as één periode van de \omterm{def:b2:orbital-forcing:elements}{klimatologische precessie} geven; het gemiddelde is de
\num{23000} jaar die gewoonlijk wordt opgegeven.
\end{solution}

\begin{solution}{exo:b2:orbital-forcing:4}
$\sqrt{1 - 0.0167^2} = 0.99986$. Voor \ang{22.1}: $1360 \times 0.37622/(\pi \times 0.99986) =
\qty{162.9}{W.m^{-2}}$; voor \ang{24.5}: $1360 \times 0.41469/(\pi \times 0.99986) =
\qty{179.5}{W.m^{-2}}$. De verhouding is $\sin\ang{24.5}/\sin\ang{22.1} = 1.102$: het jaargemiddelde op de
pool groeit met \qty{10.2}{\percent}, ongeveer \qty{16.7}{W.m^{-2}}, tussen de twee uitersten.
\end{solution}

\begin{solution}{exo:b2:orbital-forcing:5}
(a) Bij de junizonnewende is $(a/r)^2 = \bigl((1 - e\sin\varpi)/(1 - e^2)\bigr)^2$. Vandaag is $e\sin\varpi
= 0.016898$ en $(0.983102/0.999703)^2 = 0.9671$, tegenover $1 - 0.033796 = 0.9662$;
\num{116000} jaar geleden $e\sin\varpi = 0.043939$ en $(0.956061/0.998065)^2 = 0.9176$, tegenover $1 -
0.087878 = 0.9121$. (b) Bij de decemberzonnewende is $(a/r)^2 = \bigl((1 + e\sin\varpi)/(1 -
e^2)\bigr)^2$: vandaag $1.0347$ tegenover $1.0338$, \num{116000} jaar geleden $1.0940$ tegenover $1.0879$.
(c) De verwaarloosde termen zijn van orde $e^2$: $(1 - x)^2/(1 - e^2)^2 \approx 1 - 2x + x^2 + 2e^2$
met $x = e\sin\varpi$, hier ongeveer $3e^2$. Ze bedragen $3 \times 0.0172^2 = 0.0009$ vandaag en $3
\times 0.0440^2 = 0.0058$ \num{116000} jaar geleden, dicht bij de gevonden verschillen, $0.0009$ en
$0.0055$: de benadering is zes keer slechter wanneer de excentriciteit tweeënhalf keer
groter is.
\end{solution}

\begin{solution}{exo:b2:orbital-forcing:6}
Bij de junizonnewende is $\delta = \varepsilon = \ang{23.446}$. $\cos h_0 = -\tan\ang{65}\tan\ang{23.446}
= -2.14451 \times 0.43369 = -0.93006$, dus $h_0 = \ang{158.44} = \qty{2.7654}{rad}$. De dag duurt
$2 \times 158.44/15 = 21.13$~uur (de aarde draait \ang{15} per uur). Met $\sin\ang{65} = 0.90631$,
$\cos\ang{65} = 0.42262$, $\sin\delta = 0.39788$, $\cos\delta = 0.91744$ en $\sin h_0 = 0.36742$
is de haakjesterm $2.7654 \times 0.90631 \times 0.39788 + 0.42262 \times 0.91744 \times 0.36742 = 0.99720
+ 0.14246 = 1.13966$. Op de gemiddelde afstand is $\bar Q = (1360/\pi) \times 1.13966 = 432.90 \times 1.13966
= \qty{493.4}{W.m^{-2}}$; met de afstandsfactor van \cref{exo:b2:orbital-forcing:5} is $493.4 \times
0.9671 = \qty{477.1}{W.m^{-2}}$. Ter controle geven dezelfde stappen op \ang{60}~NB \qty{475.96}{W.m^{-2}}, waar de tabellen \num{475.95} geven.
\end{solution}

\begin{solution}{exo:b2:orbital-forcing:7}
(a) De maartequinox is $\lambda = 0$, de septemberequinox $\lambda = \pi$. Met $\sqrt{1 - e^2}
\approx 1$ is de duur $\frac{T}{2\pi}\int_0^{\pi}\bigl(1 + 2e\cos(\lambda -
\varpi)\bigr)\,\mathrm d\lambda = \frac{T}{2\pi}\bigl(\pi + 2e[\sin(\lambda - \varpi)]_0^{\pi}\bigr)
= \frac{T}{2\pi}(\pi + 4e\sin\varpi) = \frac T2 + \frac{2T}{\pi}e\sin\varpi$. (b) Het andere halfjaar
is evenveel korter, dus het verschil is $(4T/\pi)\,e\sin\varpi = (4 \times
365.2422/\pi) \times 0.016898 = 465.04 \times 0.016898 = 7.86$~dagen. (c) Lente en zomer samen
zijn $4.5 + 3 = 7.5$~dagen langer dan herfst en winter; de twee opgegeven verschillen zijn afgerond, en de
overeenstemming valt binnen die afronding.
\end{solution}

\begin{solution}{exo:b2:orbital-forcing:8}
(a) Voor $e = 0.0046$: $\sqrt{1 - e^2} = 0.999989$ en $340/0.999989 = \qty{340.004}{W.m^{-2}}$. Voor
$e = 0.0565$: $\sqrt{1 - 0.003192} = 0.998403$ en $340/0.998403 = \qty{340.544}{W.m^{-2}}$. (b) Het
verschil is \qty{0.54}{W.m^{-2}}, \qty{0.16}{\percent}. (c) De instraling van half juli op \ang{65}~NB
schommelde met $497 - 388 = \qty{109}{W.m^{-2}}$, ongeveer \qty{25}{\percent} van haar gemiddelde en tweehonderd
keer de verandering van het globaal jaarlijks gemiddelde: de baan werkt op het klimaat door zonlicht te verplaatsen tussen
seizoenen en breedten, niet door het totaal te veranderen.
\end{solution}

\begin{solution}{exo:b2:orbital-forcing:9}
Met $\omega = 2\pi/P$ en $\tau = \qty{10}{kyr}$: voor $P = 23$ is $\omega\tau = 2.732$, een relatieve
amplitude $1/\sqrt{1 + 7.464} = 0.344$ en een vertraging $\arctan(2.732)/\omega = 1.2199/0.2732 = 4.47$~kyr;
voor $P = 41$ is $\omega\tau = 1.532$, amplitude $0.546$, vertraging $0.9926/0.1532 = 6.48$~kyr; voor $P =
100$ is $\omega\tau = 0.628$, amplitude $0.847$, vertraging $0.5610/0.0628 = 8.93$~kyr. De ijskap
begunstigt de trage forcering: ze volgt \qty{85}{\percent} van een cyclus van \num{100000} jaar maar slechts een
derde van een precessiecyclus, omdat ze te weinig tijd heeft om te groeien of te krimpen voordat een snelle forcering
omkeert. Ze werkt als een laagdoorlaatfilter.
\end{solution}

\begin{solution}{exo:b2:orbital-forcing:10}
(a) De vergelijking is lineair: als $V_1$ antwoordt op $A_1\sin\omega_1t$ en $V_2$ op
$A_2\sin\omega_2t$, antwoordt $V_1 + V_2$ op hun som. Volgens \cref{prop:b2:orbital-forcing:lag} zijn $V_1$ en
$V_2$ sinusoïden met frequenties $\omega_1$ en $\omega_2$ (zodra het overgangsverschijnsel is uitgestorven), dus
de respons bevat geen andere frequentie. (b) Met $V \approx B_1\sin(\omega_1t - \phi_1) +
B_2\sin(\omega_2t - \phi_2)$ is
\[
\begin{aligned}
  V^2 &= \tfrac12(B_1^2 + B_2^2) - \tfrac12B_1^2\cos(2\omega_1t - 2\phi_1) -
  \tfrac12B_2^2\cos(2\omega_2t - 2\phi_2)\\
  &\quad + B_1B_2\bigl[\cos\bigl((\omega_1 - \omega_2)t - \phi_1 + \phi_2\bigr) -
  \cos\bigl((\omega_1 + \omega_2)t - \phi_1 - \phi_2\bigr)\bigr],
\end{aligned}
\]
dus de term $kV^2$ forceert de vergelijking bij de verschil- en somfrequenties, en bij de dubbele
frequenties. (c) $1/(1/19.0 - 1/23.7) = 1/(0.05263 - 0.04219) = 95.8$ duizend jaar; het paar 19.0
en 22.4 geeft 125 duizend jaar. De verschilfrequenties van de precessietermen vallen in de
excentriciteitsband, zoals het moet, want de \omterm{def:b2:orbital-forcing:elements}{klimatologische precessie} is $e\sin\varpi$. Een niet-lineaire
ijskap zou dus een ritme van \num{100000} jaar uit de precessie kunnen halen; een lineaire
niet. De cycli van \num{100000} jaar wijzen op niet-lineair gedrag, al toont deze berekening
alleen niet welk.
\end{solution}

\begin{solution}{exo:b2:orbital-forcing:11}
(a) Zoals in het bewijs van \cref{prop:b2:orbital-forcing:annual} is $(a/r)^2\,\mathrm dt =
T\,\mathrm d\lambda/(2\pi\sqrt{1 - e^2})$, en $\sin\delta = \sin\varepsilon\sin\lambda$ is positief
voor $0 < \lambda < \pi$. Dus $E = \frac{S_0T\sin\varepsilon}{2\pi\sqrt{1 - e^2}}\int_0^{\pi}\sin\lambda
\,\mathrm d\lambda = S_0T\sin\varepsilon/(\pi\sqrt{1 - e^2})$. (b) $T = 365.25 \times 86400 =
\qty{3.1558e7}{s}$, $S_0T/\pi = \qty{1.3661e10}{J.m^{-2}}$; voor \ang{22.1} is $E = 1.3661 \times 10^{10}
\times 0.37622/0.99986 = \qty{5.14}{GJ.m^{-2}}$; voor \ang{24.5} $\qty{5.67}{GJ.m^{-2}}$,
\qty{10}{\percent} meer. (c) De precessie bepaalt hoe dezelfde energie over de tijd wordt gespreid: een zomer
in het perihelium is korter en intenser, een zomer in het aphelium langer en zwakker. Smelten begint
pas boven \qty{0}{\celsius}, dus een korte intense zomer kan meer doen smelten dan een lange koele die
hetzelfde totaal ontvangt.
\end{solution}

\begin{solution}{exo:b2:orbital-forcing:12}
(a) $3.93\log_2(279.4/189.7) = 3.93\log_2 1.4729 = 3.93 \times 0.5586 = \qty{2.20}{W.m^{-2}}$. (b) De
kleinste stijging is van 194 naar \qty{266}{ppm}: $3.93\log_2 1.3711 = 3.93 \times 0.4554 =
\qty{1.79}{W.m^{-2}}$; de grootste van 188 naar \qty{282}{ppm}: $3.93\log_2 1.5 = 3.93 \times 0.5850 =
\qty{2.30}{W.m^{-2}}$. (c) De forcering van koolstofdioxide is vier keer de grootste verandering van de
globale jaarlijkse instraling, \qty{0.54}{W.m^{-2}}, en ze werkt overal en het hele jaar; de zomer%
instraling op \ang{65}~NB schommelt met ongeveer \qty{109}{W.m^{-2}}, maar op één breedte en in één seizoen. De
baan is de trekker, via de zomers van de noordelijke ijskappen; koolstofdioxide is, met de
\omterm{def:b2:precambrian-earth:ice-albedo}{ijs--albedoterugkoppeling}, een versterker die de verandering over de aarde verspreidt.
\end{solution}

\begin{solution}{pb:b2:orbital-forcing:1}
\textbf{1.} Toen $0.0440 \times \sin\ang{92.7} = 0.0440 \times 0.99889 = 0.04395$; nu $0.0172 \times
\sin\ang{101.4} = 0.0172 \times 0.98027 = 0.01686$, ongeveer 2.6 keer minder.
\textbf{2.} Toen $1 - 0.0440 = 0.956$ en $1 + 0.0440 = 1.044$; nu $0.9828$ en $1.0172$.
\textbf{3.} Bij $\lambda = \varpi + \ang{180}$: toen \ang{272.7}, nu \ang{281.4}, beide dicht bij de
decemberzonnewende ($\lambda = \ang{270}$); toen passeerde de aarde het perihelium minder dan drie dagen
erna, en het aphelium bij de junizonnewende.
\textbf{4.} Toen $(1.044/0.956)^2 = 1.0921^2 = 1.193$, nu $(1.0172/0.9828)^2 = 1.0350^2 = 1.071$:
toen \qty{19}{\percent} meer zonlicht in het perihelium, nu \qty{7}{\percent}.
\textbf{5.} Kleiner, \ang{22.52} tegenover \ang{23.45}: de zomerzon stond lager op hoge breedten,
en hoge breedten ontvingen minder zonlicht in de zomer en over het jaar.
\textbf{6.} Toen $340/\sqrt{1 - 0.001936} = 340/0.999032 = \qty{340.33}{W.m^{-2}}$; nu $340/0.999852
= \qty{340.05}{W.m^{-2}}$. Het verschil, \qty{0.28}{W.m^{-2}} of minder dan een duizendste,
is onbeduidend.
\textbf{7.} Bij $\lambda = \ang{90}$ is $\cos(\ang{90} - \varpi) = \sin\varpi$ in
\cref{prop:b2:orbital-forcing:distance}.
\textbf{8.} Toen $(1 - 0.001936)/(1 - 0.04395) = 0.998064/0.956049 = 1.04395$; nu $(1 -
0.000296)/(1 - 0.01686) = 0.999704/0.983139 = 1.01685$.
\textbf{9.} Toen $1/1.04395^2 = 0.91758$; nu $1/1.01685^2 = 0.96714$.
\textbf{10.} Toen $1 - 2 \times 0.04395 = 0.91210$; nu $1 - 2 \times 0.01686 = 0.96628$: de
benadering wijkt toen $0.0055$ af, nu $0.0009$.
\textbf{11.} Toen $1 - 0.91758 = \qty{8.2}{\percent}$ zwakker, nu \qty{3.3}{\percent}.
\textbf{12.} Bij de zonnewende is $\delta = \varepsilon$: toen \ang{22.52}, nu \ang{23.45}.
\textbf{13.} Toen $\cos h_0 = -2.14451 \times 0.41462 = -0.88916$, $h_0 = \ang{152.77} =
\qty{2.6663}{rad}$; nu $\cos h_0 = -2.14451 \times 0.43378 = -0.93023$, $h_0 = \ang{158.47} =
\qty{2.7659}{rad}$.
\textbf{14.} $2h_0/\ang{15}$ per uur: toen $20.37$~uur, nu $21.13$~uur.
\textbf{15.} Toen $2.6663 \times 0.90631 \times 0.38301 + 0.42262 \times 0.92375 \times 0.45759 =
0.92554 + 0.17864 = 1.10418$ (\num{1.10417} onafgerond), en $432.90 \times 1.10417 =
\qty{478.0}{W.m^{-2}}$; nu $2.7659 \times 0.90631 \times 0.39795 + 0.42262 \times 0.91741 \times
0.36697 = 0.99757 + 0.14228 = 1.13985$ (\num{1.13982} onafgerond), en $432.90 \times 1.13982 = \qty{493.4}{W.m^{-2}}$.
\textbf{16.} Toen $478.0 \times 0.91758 = \qty{438.6}{W.m^{-2}}$; nu $493.4 \times 0.96714 =
\qty{477.2}{W.m^{-2}}$, \qty{38.6}{W.m^{-2}} meer dan toen.
\textbf{17.} Half juli is de lengte van de zon \ang{120}: $\sin\delta = \sin\varepsilon\sin\ang{120}$,
een declinatie van ongeveer \ang{20} in plaats van \ang{23}, dus een kortere dag en een lagere zon: op
\ang{65}~NB is de daggemiddelde instraling het grootst rond de zonnewende, wanneer de declinatie het grootst is.
\textbf{18.} Met de huidige afstand kost de verandering van de scheefstand alleen $(493.4 - 478.0) \times
0.96714 = \qty{14.9}{W.m^{-2}}$; met de oude scheefstand kost de afstand alleen $478.0 \times
(0.96714 - 0.91758) = \qty{23.7}{W.m^{-2}}$. Samen geven ze \qty{38.6}{W.m^{-2}}; de afstand
telde voor ongeveer drie vijfde.
\textbf{19.} Na een plotselinge verandering heeft de ijskap $1 - \mathrm e^{-t/\tau}$ ervan bereikt:
$1 - \mathrm e^{-0.5} = 0.39$ na \qty{5}{kyr}, $1 - \mathrm e^{-1} = 0.63$ na \qty{10}{kyr},
$1 - \mathrm e^{-2} = 0.86$ na \qty{20}{kyr}.
\textbf{20.} $\omega\tau = 2\pi \times 10/23 = 2.732$: de ijskap volgt $1/\sqrt{1 + 2.732^2} =
0.34$ van de forcering, $\arctan(2.732)/\omega = 4.5$~kyr te laat.
\textbf{21.} Een halve precessiecyclus, zo'n 11 tot \qty{12}{kyr}, is ongeveer één responstijd: voordat
de ijskap twee derde van de weg heeft afgelegd, keert de instraling om en trekt haar terug.
\textbf{22.} In het lineaire model groeit en krimpt het ijs even snel bij forceringen van dezelfde
grootte: een stijgende en een dalende forcering geven spiegelbeeldige responsen. Een trage opbouw gevolgd door een snelle
ineenstorting vraagt een respons die afhangt van de toestand of van de richting van de verandering, dus een
niet-lineaire.
\textbf{23.} Koolstofdioxide steeg met $3.93\log_2(266/194) = 1.8$ tot $3.93\log_2(282/188) =
\qty{2.3}{W.m^{-2}}$ van glaciaal naar \omterm{def:b2:orbital-forcing:glacials}{interglaciaal}: een forcering over de hele aarde en het hele
jaar. Vrijgekomen wanneer de Atlantische circulatie verzwakt, draagt het de verandering van de noordelijke
zomers over naar de hele planeet, en versterkt het die, samen met het albedo van het krimpende ijs.
\textbf{24.} Toen het laatste \omterm{def:b2:orbital-forcing:glacials}{interglaciaal} eindigde, bedroeg de daggemiddelde instraling bij de junizonnewende op
\ang{65}~NB ongeveer \textbf{\qty{439}{W.m^{-2}}}, \qty{39}{W.m^{-2}} minder dan nu en haar laagste
waarde van de cyclus.
\end{solution}
